%\ifodd\value{page}\hbox{}\newpage\fi
\chapter*{Śrī Lalitā Sahasranāma — Śloka 6}
\addcontentsline{toc}{chapter}{Śloka 6 - Nomi 14,15}

\begin{sloka}
{\sanskritfont
\begin{tabular}{c}
वदनस्मरमाङ्गल्यगृहतोरणचिल्लिका ।\\
वक्त्रलक्ष्मीपरीवाहचलन्मीनाभलोचना ॥ ६ ॥
\end{tabular}

\vspace{1.5em}

{\middlesize \iastfont
vadana-smara-māṅgalya-gṛha-toraṇa-cillikā |\\
vaktra-lakṣmī-parīvāha-calan-mīnābhā-locanā ||6||}

\vspace{1em}

{\middlesize
\anvaya{%
वदनस्मरमाङ्गल्यगृहतोरणचिल्लिका ।
वक्त्रलक्ष्मीपरीवाहचलन्मीनाभलोचना ।
}
}

\middlesize \iastfont \itmean{%
«Il cui volto è arco augurale all’ingresso della dimora dell’amore;  
i cui occhi, simili a pesci splendenti, si muovono nel fluire dell’abbondanza del volto.»%
}

}
\end{sloka}

\wordmeaning{%
\wordline{वदन-स्मर-माङ्गल्य-गृह-तोरण-चिल्लिका}
{vadana-smara-māṅgalya-gṛha-toraṇa-cillikā}
{ il cui volto (\textit{vadana}) è come l’arco ornamentale (\textit{toraṇa-cillikā}) della dimora (\textit{gṛha}) propizia (\textit{māṅgalya}) di \textit{Smara}, il desiderio divinizzato; }%

\wordline{वक्त्र-लक्ष्मी-परिवाह-चलन्-मीन-आभा-लोचना}
{vaktra-lakṣmī-parīvāha-calan-mīnābhā-locanā}
{ i cui occhi (\textit{locanā}) possiedono lo splendore (\textit{ābhā}) dei pesci (\textit{mīna}), muovendosi (\textit{calan}) nel flusso (\textit{parīvāha}) della bellezza (\textit{lakṣmī}) del volto (\textit{vaktra}); }%
}



Lo \textit{Śloka} 6 approfondisce la manifestazione della bellezza di \textit{Lalitā} come principio ordinatore e fecondante, non come ornamento estetico isolato. L’espressione \textit{vadana-smara-māṅgalya-gṛha-toraṇa-cillikā} costruisce un’immagine architettonica precisa: il volto non è semplice bellezza, ma soglia, punto di passaggio, luogo in cui l’energia del desiderio (\textit{smara}) viene resa propizia, ritmata e ordinata.

La seconda metà dello \textit{śloka} introduce il movimento interno di tale bellezza. L’espressione \textit{vaktra-lakṣmī-parīvāha} può essere intesa come il fluire dell’abbondanza del volto, la corrente della \textit{śrī}, non qualità statica ma presenza dinamica che permea e sostiene l’intera fisionomia di \textit{Lalitā}. Gli occhi, descritti come \textit{calan-mīnābhā}, partecipano a questo flusso come pesci luminosi immersi nella corrente stessa della \textit{śrī-lakṣmī}.

Il simbolismo dei \textit{mīna} non indica soltanto mobilità o grazia visiva, ma richiama anche il respiro vitale duale che anima la percezione: i due flussi complementari del prāṇa, o, in linguaggio yogico, le correnti di \textit{iḍā} e \textit{piṅgalā}. Il volto, sede della parola e dello sguardo, diviene così il luogo in cui il ritmo vitale della coscienza si rende percepibile, senza separazione tra energia, forma e senso.

In questo \textit{śloka}, la bellezza di \textit{Lalitā} non è decorativa ma ritmica: essa manifesta l’ordine cosmico del \textit{śṛṅgāra-rasa} divino, in cui attrazione, respiro e percezione sono modulazioni della stessa \textit{Cit-śakti}. La bellezza è dunque funzione attiva della coscienza: non semplice atto volontario, ma armonia vivente che regge e anima il mondo.

Questo \textit{śloka} contiene due \textit{nāma}: \textbf{\textit{Vadana-smara-māṅgalya-gṛha-toraṇa-cillikā}} definisce il volto come soglia augurale della dimora dell’amore, indicando che il desiderio viene assunto e ordinato in forma propizia; \textbf{\textit{Vaktra-lakṣmī-parīvāha-calan-mīnābhā-locanā}} descrive gli occhi come movimento interno della \textit{śrī-lakṣmī} del volto: non semplice grazia iconografica, ma ritmo percettivo e vitale che scorre nella fisionomia divina.
